\scanpage{112}
\mspage{55}
compatible avec les opérations de $\unsure{\mathfrak S|\mathcal S}_d$, d'où un morphisme
\[
(162)\quad\diagram{\begin{tikzcd}[ampersand replacement=\&, column sep=large,row sep=small]
D_J\arrow[r,"f_x"]\arrow[d,phantom,"\wr" description]\&D_{d,\delta}\deleted{\unsure{'}}\arrow[d,phantom,"\wr" description]\\
\widetilde D_d\wedge_{\unsure{\mathfrak S|\mathcal S}_d}H(J)\&\unsure{\widetilde D'_d|\widetilde{\widetilde D}_d}\wedge_{\unsure{\mathfrak S|\mathcal S}_d}\mathcal P_\delta(I_d)
\end{tikzcd}}
\]
\deleted{Ceci posé, il} Cette application est décrite entièrement en termes du rev. principal $\widetilde D_d$ de $D_d$, de groupe $\unsure{\mathfrak S|\mathcal S}_d$, \inserted{et de $x\in J$.} On a d'autre part, par ``fonctorialité \inserted{en $J$ de $D_J$}'', les \deleted{morphismes}\inserted{diagrammes} commutatifs étendant (158)
\[
\deleted{(\unsure{162})\qquad D_{d,\unsure{d'|\delta}}\longrightarrow D_{d'}}
\]
\deleted{et on a}
\[
(163)\quad\diagram{\begin{tikzcd}[ampersand replacement=\&]
D_J\arrow[r]\arrow[d]\&D_{J'}\arrow[d,"f_x"]\\
D_{d,\delta}\arrow[r]\&D_\delta
\end{tikzcd}},
\]
où $D_J-D_{d,\delta}$ et $D_{d,\delta}-D_\delta$ sont définis par fonctorialité par les inclusions $\deleted{\unsure{\{e\}}}\{x\}\subset\{x,e\}\subset J$. On \deleted{Cela signifie}\inserted{\unsure{peut voir aussitôt}} que $D_J\unsure{-|\to}D_{d,\delta}$ est décrit par la \unsure{nôtre|flèche} (162), donc le composé (158) est connu quand on connaît les $D_{d,\delta}\longrightarrow D_\delta$
\margin{111}
